\section{Atomi a più elettroni}
In questi atomi bisogna tenere conto delle possibili interazioni di spin-spin, spin-orbita e dei momenti angolari orbitali tra gli elettroni. Queste interazione le indichiamo operatorialmente come 
\begin{equation*}
  c_{ik} \bm{s_i}\cdot \bm{s_k} \quad a_{ij} \bm{l}_i \cdot\bm{s}_j \quad b_{ij} \bm{l}_i\cdot \bm{l}_j
\end{equation*}
L'obiettivo, allora, è quello di individuare i livelli di energia tenuto conto dell' effetto di ciascun termine. Nel fare questo possiamo trovare due situazione estreme:
\begin{itemize}
  \item Il primo regime è quello in cui l' interazione spin-spin è molto più forte dell' interazione spin-orbita $|c_{ik}|\gg |a_{ik}|$, in questo caso si può trascurare l' interazione spin-orbita e considerare solo gli effetti di spin-spin e orbita-orbita. Possiamo in questo regime definire gli operatori 
    \begin{equation*}
      \bm{L} = \sum \bm{l}_i \quad \bm{S} = \sum \bm{s}_i
    \end{equation*}
    Andiamo a definire una Hamiltoniana totale, efficace, che descrive l' interazione Spin-Orbita totale 
    \begin{equation*}
      H_{SO} = \xi_{LS} \bm{L}\cdot \bm{S}
    \end{equation*}
    Per diagonalizzare questa Hamiltoniana non vanno più bene gli stati $\ket{L,S;M_L;M_S}$, ma devo classificare gli stati in una nuova base definita dal momento angolare totale
    \begin{equation*}
      \bm{J} = \bm{S}+\bm{L}
    \end{equation*}
    $\left\{\ket{L,S;J,M_J}\right\}$. In questo regime vale lo schema LS.
  \item Il secondo regime invece vale che l' interazione di spin-orbita di singolo elettrone è non trascurabile e molto più grande delle altre interazioni
    \begin{equation*}
      |a_{ik}|\gg |c_{ik}|
    \end{equation*}
    In questo regime allora dobbiamo considerare il momento angolare totale di ciascun elettrone
    \begin{equation*}
      \bm{j}_i = \bm{s}_i + \bm{l}_i
    \end{equation*}
    In questo regime si accoppiano i momento angolari totali dei singoli elettroni e quindi l' Hamiltoniana efficace risulta
    \begin{equation*}
      H
    \end{equation*}
      Anche in questo caso definiamo un momento angolare totale, costruito diversamente dal caso dello schema LS 
      \begin{equation*}
        \bm{J} = \sum \bm{j}_i
      \end{equation*}
      Gli stati si classificano come 
      \begin{equation*}
        (j_1,\dots,j_N)_J
      \end{equation*}
\end{itemize}
